% ============================================================================ %
%     プリアンブル
% ============================================================================ %
\documentclass[
  10pt,       % フォントサイズ
  twocolumn   % ２段組（２カラム）
]{jlreq}      % 使用するフォント？


% ---------------------------------------------------------------------------- %
%     「定義」などの環境の定義
% ---------------------------------------------------------------------------- %
\newtheorem{dfn}{定義}
\newtheorem{lem}{補題}
\newtheorem{thrm}{定理}
\newtheorem{pblm}{問題}
\newtheorem{exmp}{例}


% ---------------------------------------------------------------------------- %
%     数学系の環境の定義
% ---------------------------------------------------------------------------- %
\usepackage{amsmath}
\usepackage{amssymb}
\usepackage{enumitem}
\usepackage{latexsym}
\usepackage{mathtools}
\usepackage{breqn}

\DeclarePairedDelimiter{\abs}{\lvert}{\rvert}      % 例) \abs{X}            -> |X|
                                                   %     \abs[\Big]{X}      -> big(|) X big(|)
                                                   %     \abs*{\frac{a}{b}} -> |a/b|
\DeclarePairedDelimiter{\set}{\lbrace}{\rbrace}    % 例) \set{a, b}  -> {a, b}
\newcommand{\true}{\textit{true}}
\newcommand{\false}{\textit{false}}
\newcommand{\argmax}{\mathop{\rm arg~max}\limits}
\newcommand{\argmin}{\mathop{\rm arg~min}\limits}


% ---------------------------------------------------------------------------- %
%     疑似コード関連のコマンド定義
%       参考サイト
%         https://jdhao.github.io/2019/09/21/latex_algorithm_pseudo_code/
% ---------------------------------------------------------------------------- %
\usepackage[linesnumbered, longend, ruled, vlined]{algorithm2e}

\SetKwInOut{KwIn}{Input}
\SetKwInOut{KwOut}{Output}
\newcommand{\Break}{break}
\newcommand{\Continue}{continue}
\newcommand{\Call}[2]{\textup{#1}(#2)}          % 例) \Call{hoge}{x, y} -> hoge(x, y)
\newcommand{\Function}[3]{%                       例) \Function{hoge}{x}{body}
  \SetKwProg{Fn}{define}{}{end define}%
  \Fn{\textup{#1(#2)}}{#3}%
}
\newcommand{\Goto}{%
  \SetKw{KwGoto}{goto}%
  \KwGoto%
}
\newcommand{\BMatrix}[1]{%
  $\begin{bmatrix}%
    #1%
  \end{bmatrix}$%
}
\newcommand{\Mfunc}[1]{%
  \textsc{M}(#1)%
}
\renewcommand{\algorithmcfname}{アルゴリズム} % Algorithmをアルゴリズムにする（英語を無理やり日本語にする）．


% ---------------------------------------------------------------------------- %
%     キャプションと参照コマンドの定義
% ---------------------------------------------------------------------------- %
\newcommand{\secref}[1]{\ref{#1}章}
\newcommand{\subsecref}[1]{\ref{#1}節}
\newcommand{\subsubsecref}[1]{\ref{#1}項}
\newcommand{\figref}[1]{図\ref{#1}} % 図は下に...
\newcommand{\tblref}[1]{表\ref{#1}} % 表は上に...
\newcommand{\algref}[1]{アルゴリズム\ref{#1}}
\newcommand{\coderef}[1]{\ref{#1}行}
\newcommand{\dfnref}[1]{定義\ref{#1}}
\newcommand{\lemref}[1]{補題\ref{#1}}
\newcommand{\thrmref}[1]{定理\ref{#1}}
\newcommand{\pblmref}[1]{問題\ref{#1}}
\newcommand{\exmpref}[1]{例\ref{#1}}


% ---------------------------------------------------------------------------- %
%     余白設定
%       参考サイト
%         https://cns-guide.sfc.keio.ac.jp/2001/11/5/1.html
% ---------------------------------------------------------------------------- %
\setlength{\topmargin}{-2.5cm}  % 各ページの上側の余白．
\setlength{\textheight}{27cm}
\setlength{\textwidth}{18.5cm}
\setlength{\oddsidemargin}{-1.3cm}
\setlength{\columnsep}{.5cm}
\catcode`@=\active \def@{\hspace{0.9bp}-\hspace{0.9bp}}


% ---------------------------------------------------------------------------- %
\setlength{\topmargin}{-2.5cm}  % 各ページの上側の余白．
\setlength{\textheight}{27cm}
\setlength{\textwidth}{18.5cm}
\setlength{\oddsidemargin}{-1.3cm}
\setlength{\columnsep}{.5cm}
\catcode`@=\active \def@{\hspace{0.9bp}-\hspace{0.9bp}}


% ---------------------------------------------------------------------------- %
%     タイトル
% ---------------------------------------------------------------------------- %
\title{
  研究環境のクラウド化 \\
}
\author{
 田中研究室：佐々木ギリェルミ
}
\date{
  %\today  % \todayを書くことでコンパイルした日付が自動入力されます．
}


% ============================================================================ %
%     本文作成
% ============================================================================ %
\begin{document}
\maketitle  % \begin{document} ~ \end{document}内に\maketitleを書くことで，タイトルがページトップに自動作成されます．

\section{はじめに}\label{sec:Section}  % \label{<任意の文字列>}でラベル付けできます．後から参照するときに自動で番号を振ってくれるので便利です．\labelは後から参照したいものの直後に付けます．
現在、本研究室では、研究活動に用いる文書、ソースコード、実験データ、発表資料および開発環境に関する情報といった研究室リソースを、研究室内に設置されたサーバーを通して共有している。研究室内サーバーを利用することで、研究室リソースを一か所に集約できる一方、原則として研究室内からの利用を前提としているため、研究室外から資料を確認することや、研究室のメンバー以外と情報を共有することが容易ではないという問題があった。

そして、本研究室では、学生が毎年それぞれ独立した研究を新たに開始するのではなく、複数年度にわたる研究プロジェクトを引き継ぎながら継続して進めている。そのため、卒業や修了によって担当者が交代する際には、研究の目的や背景だけでなく、これまでに作成された資料、ソースコード、実験結果および開発環境の構築方法などを、次の担当者へ正確に引き継ぐ必要がある。

そこで本取り組みでは、研究室リソースを研究プロジェクトごとに整理し、GitHub上で管理・共有できる環境を構築することを目的とした。これにより、研究資料の保存場所を明確にし、研究室外からの情報共有を容易にするとともに、研究内容および開発環境の円滑な引き継ぎを実現することを目指した。

\section{既存環境の課題}\label{sec:Section}  % \label{<任意の文字列>}でラベル付けできます．後から参照するときに自動で番号を振ってくれるので便利です．\labelは後から参照したいものの直後に付けます．
研究室内サーバーに保存されているファイルを確認した結果、研究プロジェクトと保存されているファイルの対応関係が分かりにくいという課題が存在した。

フォルダ名だけでは内容を判断できない場合があるほか、サーバー上で使用されているユーザーIDと学生名が一致しないため、どの学生がどの研究を担当していたのかを特定することが困難であった。そのため、必要な資料を探すために多くの時間を要していた。

さらに、長期間継続している研究プロジェクトでは、歴代の担当者が作成した資料やプログラムが蓄積されており、現在の担当者が前任者の研究内容を理解し、引き継ぐことが大きな負担となっている。

以上の課題から、研究資料を研究プロジェクトごとに分類し、ファイルの保存場所と管理方法を統一するとともに、担当者が交代した後も継続的に利用できる管理環境が必要であると考えた。

\section{管理環境の設計}\label{sec:Section}  % \label{<任意の文字列>}でラベル付けできます．後から参照するときに自動で番号を振ってくれるので便利です．\labelは後から参照したいものの直後に付けます．
本取り組みでは、研究資料の管理基盤としてGitHubを採用した[1]。GitHubでは、研究内容ごとにリポジトリを作成できるほか、Gitを利用してファイルの変更履歴を記録できる。そのため、ファイルを更新した日時、更新者および変更内容を確認でき、従来のようにファイル名のみで版を管理する必要がない。また、インターネットを通して利用できるため、適切なアクセス権限を設定することで、研究室外からの確認や外部の共同研究者との情報共有も可能となる。

\section{GitHubの概要}\label{sec:Section}  % \label{<任意の文字列>}でラベル付けできます．後から参照するときに自動で番号を振ってくれるので便利です．\labelは後から参照したいものの直後に付けます．
GitHubは、ソースコードやファイルをクラウド上で管理・共有するためのサービスである。データの変更履歴を保存できるほか、複数人で同じプロジェクトを管理できるため、外部研究者との情報共有や学生間の引継ぎを円滑に行うことができる。

\section{資料の分類}\label{sec:Section}  % \label{<任意の文字列>}でラベル付けできます．後から参照するときに自動で番号を振ってくれるので便利です．\labelは後から参照したいものの直後に付けます．
はじめに、研究室内サーバーに保存されているファイルの内容を確認し、どの研究プロジェクトに関係するファイルであるかを分類した。分類の際には、フォルダ名やファイル名だけでなく、ファイルの内容、作成者および更新日時なども確認した。

次に、調査した研究内容を分析し、関連するテーマごとに分類を行った。その結果、研究テーマを17のグループに整理することができた。

最後に、分類した各研究テーマについてGitHub上にリポジトリを作成し、研究資料やソースコードを保存・管理できる環境を構築した。これにより、研究データの一元管理が可能となり、外部研究者との共有や学生間の円滑な引継ぎを行える基盤を整備した。

\section{まとめと今後の課題}\label{sec:Section}  % \label{<任意の文字列>}でラベル付けできます．後から参照するときに自動で番号を振ってくれるので便利です．\labelは後から参照したいものの直後に付けます．
本研究では、研究室内サーバーに保存されていたファイルを調査し、研究プロジェクトごとに分類するとともに、研究テーマごとにGitHubリポジトリを作成した。これにより、各研究に関する資料の保存場所が明確になり、ファイルの検索性、共有性および変更履歴の管理性の向上が期待できる。また、研究テーマごとに整理することで、研究内容の把握や学生間の引継ぎが容易になる環境を整備した。

今後の課題として、GitHub上に作成した各リポジトリへ研究資料やソースコードなどの研究リソースを移行する作業が挙げられる。また、リポジトリの運用ルールやディレクトリ構成を整備し、継続的に利用・管理できる環境を構築する必要がある。さらに、外部研究者との共同利用を想定し、アクセス権限や共有方法についても検討を進める予定である。

\subsection{分類}\label{sec:Section}  % \label{<任意の文字列>}でラベル付けできます．後から参照するときに自動で番号を振ってくれるので便利です．\labelは後から参照したいものの直後に付けます．
1. ルール最適化・アルゴリズム	フィルタリングルールの並べ替えや統合など、ルール数や処理効率を改善するアルゴリズムに関する研究。
2. ルール順序最適化	ルールの順番を変更することで検索時間や処理速度を向上させる研究。
3. ルールリスト最適化	ルールの生成・分割・統合など、ルールリスト自体を改善する研究。
4. フレームワーク	フィルタリングルール最適化を支援するシステムやフレームワークの構築。
5. スケジューリング	フローやタスクの最適な割り当て・スケジューリングに関する研究。
6. 決定木	決定木やTrieなどを利用した高速なパケット分類手法の研究。
7. Run-Based Trie	Run-Based Trieを用いた検索・決定木構築・性能評価に関する研究。
8. SGM	SGMアルゴリズムの改良や重み付け、ルール選択手法に関する研究。
9. EAD	Extensional Acyclic Digraph（EAD）の構築・解析・故障検知などの研究。
10. ClassBench	ClassBenchを利用したフィルタリングルール生成や評価環境に関する研究。
11. パケット分類・空間分割	パケット分類や空間分割アルゴリズムによる高速検索手法の研究。
12. テスト環境・実証実験	ネットワーク実験環境の構築や最適化手法の評価・検証に関する研究。
13. 攻撃防御・セキュリティ応用	DDoS対策、IDS、ウイルス検知などセキュリティ技術の研究。
14. 機械学習・AI	機械学習を利用したルール最適化やソフトウェア解析に関する研究。
15. NFV・仮想ネットワーク	NFV、OpenFlow、仮想ネットワーク埋め込みなどの研究。
16. その他・補助的テーマ	上記に分類されない補助的な研究や関連技術の研究。
17. 独立テーマ	他の分類に属さない独立したテーマの研究。

\section{参考文献}\label{sec:Section}  % \label{<任意の文字列>}でラベル付けできます．後から参照するときに自動で番号を振ってくれるので便利です．\labelは後から参照したいものの直後に付けます．
[1]GitHub, Inc., GitHub, https://github.com/
   

\end{document}

